\begin{center}
\textbf{\large Abstract}
\end{center}
\addcontentsline{toc}{chapter}{Abstract}
\begin{small}

Python implementations can agree on every base test and still disagree on a new input. This thesis uses large language models to find such inputs. Equivalence partitioning and boundary value analysis share a prompt that asks for class representatives and values near their boundaries. A second approach derives inputs from the rules of a decision table. ControlFlow targets a line or branch the base tests missed and uses execution traces to revise its proposals. The HumanEval+ reference implementation supplies expected results once the input set is fixed.

Of 1,454 eligible programs across 157 tasks, the three-branch pipeline found discrepancies in 121 with Qwen-generated tests and in 132 with Gemini-generated tests. On one word-splitting task, a tab helped ControlFlow find eight implementations missed by the other two Qwen branches. Across both generators, the pipeline found 149 distinct programs and the direct-prompting baseline found 120. On the 162 tasks with ten Qwen implementations each, adding the pipeline tests lowered pass@1 from 89.75\% to 80.56\%. The contribution of each branch under a common budget remains open.

\vspace*{.6cm}
\textbf{Keywords}\quad test input generation, large language models, boundary value analysis, decision tables, branch coverage.

\end{small}
